\chapter{Compaction-dependent microstructure and hydraulic behaviours}
\label{chap:three}

\section{Introduction}
Road and rail earthworks, such as embankments and engineered slopes, form a substantial part of transport infrastructure. Their stability is governed by a climate-driven change in matric suction, hydraulic conductivity and soil-water redistribution \parencite{glendinningConstructionManagementMaintenance2014, rouainiaDeteriorationInfrastructureCutting2020}. However, recent global heatwaves and atmospheric aridity increase evaporative demand and the risk of near-surface moisture depletion \parencite{qingAcceleratingFlashDroughts2022}. These conditions create a need for lower-impact amendments that can moderate near-surface moisture loss while maintaining acceptable hydraulic performance. Within the broader move towards nature-based and bio-based geotechnical solutions \parencite{dejongBiogeochemicalProcessesGeotechnical2013}, biochar, a carbon-rich product of biomass pyrolysis, has attracted increasing attention due to its low environmental impact, carbon-sequestration potential, and capacity to modify soil hydraulic properties \parencite{sohiCarbonStorageBenefits2012, hussainInvestigatingUnsaturatedHydraulic2021, dingWaterRetentionCharacteristics2025}. Its application in geotechnical earthworks, however, requires a mechanistic understanding of how biochar dosage and compaction state modify the soil pore structure and, consequently, its \ac{SWRC} and hydraulic conductivity during drying \parencite{hussainInfluenceBiocharSoil2020}.\par

%% define the hydraulic conductivity
To investigate the coupled effects of biochar dosage and compaction state on the pore structure and soil hydraulic behaviours, this study investigates a low-plasticity sandy earthwork soil amended with 0\%, 5\% and 10\% biochar at degrees of compaction \ac{DC} of 80\% and 95\%. Drying \ac{SWRC} over an extended suction range is measured by combining continuous evaporation measurements with dew-point hygrometry, while unsaturated \ac{HC} is also determined. \Ac{MIP} and \ac{SEM} tests are used to quantify pore-domain redistribution and to identify the spatial relationships among soil aggregates, biochar particles and interfacial pores. The \ac{IMWC} and \ac{IMHC}, derived from the \ac{SWRC} and unsaturated \ac{HC} curves, are calculated to assess the changes in water-retention capacity and hydraulic conductivity. Finally, a mechanistic link between compaction, pore restructuring and hydraulic behaviours is established. In specific, three research gaps are addressed in this chapter:
\begin{itemize}
    \item The influence of initial compaction state on biochar-induced restructuring of the multiscale pore system.
     \item The pore domains governing changes in water retention and unsaturated \ac{HC} with increasing suction.
     \item The conditions under which enhanced water storage remains coupled with hydraulic transmission.
\end{itemize} 
The results will provide a mechanistic understanding of biochar application in transportation geotechnical earthworks.\par

\section{Materials and sample preparation}
\subsection{Tested soil and biochar}
The tested soil was collected from an area adjacent to the HS2 construction site near the University of Warwick, UK. Bulk soils were excavated to a depth of 1.0 m below the ground surface. After removal of visible gravels and roots, the soil was oven-dried at 105 °C for 48 h and sieved to less than 2 mm for testing. The basic properties were determined following the BS EN ISO 17892 and BS 1377 standard series. Figure \ref{fig:ch3GrainSizeDistribution}a presents its grain size distribution. Table \ref{tab:ch3soilproperties} presents the basic geotechnical properties of the tested soil. According to the BS EN ISO 14688-2, the soil exhibits a nearly medium-graded and is classified as a low plasticity SAND with silt and clay (CIL).\par
The tested biochar, provided by a commercial supplier, is derived from arboricultural arisings comprising mixed European hardwood species, including oak and robinia. It is produced through pyrolysis for 4 to 8 h in a retort kiln at around 425 °C. Its tapped bulk density and specific gravity are 380 $\mathrm{kg}/m^{3}$, and 1.07, respectively. A Bettersize 2600 laser grain size analyser is employed to determine the grain size distribution and recorded as a volume percentage. As presented in Figure \ref{fig:ch3GrainSizeDistribution}b, the grain size of biochar ranges from 1.57 $\mu$m to 1,826 $\mu$m. The incremental passing curve exhibits a bimodal pattern with two dominant particle sizes of 98.09 $\mu$m and 631.41 $\mu$m.\par

\begin{figure}[htbp]
    \centering
    \includegraphics[width=1.0\linewidth]{../figures/Ch3/ch3GrainSize.png}
    \caption{Grain size distribution of tested materials: (a) soil and (b) biochar}
    \label{fig:ch3GrainSizeDistribution}
\end{figure}
\begin{table}[htbp]
    \centering
    \caption{Basic geotechnical properties of the tested soil}
    \label{tab:ch3soilproperties}
    \begin{tabularx}{\textwidth}{>{\raggedright\arraybackslash}X >{\raggedright\arraybackslash}p{0.40\textwidth}}
        \toprule
        Parameter & Value \\
        \midrule
        Specific gravity, $G_{\mathrm{s,s}}$ & 2.73\\
        Organic matter content, $w_{\mathrm{org}}$ (\%) & 3.62 \\
        \multicolumn{2}{l}{\textit{Grading (BS EN ISO 17892-4 and BS EN ISO 14688-2)}} \\
        Sand content (\SI{63}{\micro\meter} to \SI{2}{\milli\meter}, \%) & 76.48 \\
        Silt content (\SI{2}{\micro\meter} to \SI{63}{\micro\meter}, \%) & 22.02 \\
        Clay content (\textless\SI{2}{\micro\meter}, \%) & 1.5 \\
        $D_{10}$ ($\mu$m) & 17.62 \\
        $D_{30}$ ($\mu$m) & 81.07 \\
        $D_{60}$ ($\mu$m) & 218.71 \\
        Uniformity coefficient, $C_U$ & 12.41 \\
        Coefficient of curvature, $C_C$ & 1.71 \\
        \multicolumn{2}{l}{\textit{Atterberg limits (BS 1377-2)}} \\
        Liquid limit, $w_L$ (\%) & 31.54 \\
        Plastic limit, $w_P$ (\%) & 18.79 \\
        Plasticity index, $I_P$ (\%) & 12.75 \\
        \multicolumn{2}{l}{\textit{Compaction test (BS 1377-2)}} \\
        Maximum dry density (kg/m$^3$) & 1,800 \\
        Optimum moisture content (\%) & 13.89 \\
        \bottomrule
    \end{tabularx}
\end{table}
Three biochar dosages are considered: 0\%, 5\%, and 10\%, which are the relative mass percentages of biochar to the dry soil. The specific gravities of the soil-biochar mixture are 2.73, 2.54, and 2.42 under 0\%, 5\%, and 10\% biochar dosages, respectively. Standard compaction tests were conducted to determine the \ac{MDD} and \ac{OMC} at different biochar dosages. The test results are shown in Figure \ref{fig:ch3MDDOMC}. Biochar increases the OMC and reduces the MDD, causing the optimum compaction point to shift closer to the zero-air-voids curve with increasing biochar content.\par
\begin{figure}[htbp]
    \centering
    \includegraphics[width=1.0\linewidth]{../figures/Ch3/ch3MDDOMC.png}
    \caption{Maximum dry density and optimum moisture content of biochar-amended soil at different biochar dosages: (a) pure soil, (b) 5\% biochar, and (c) 10\% biochar}
    \label{fig:ch3MDDOMC}
\end{figure}
Apparent static contact angles were measured to assess the surface wettability of the untreated soil, biochar, and biochar-amended mixtures \parencite{jingDynamicContactAngle2026}. The test was conducted on a Lauda Scientific LSA100S-T instrument (Germany). Before testing, the dried biochar and soil were ground into a fine powder and sieved through a 200-mesh screen. Approximately 200 mg of the powder was then compressed into a flat, thin pellet using a hydraulic press for 1 minute. A droplet of deionised water was deposited onto the pellet surface. Droplet images were recorded to determine the contact angle. The testing results are presented in Figure \ref{fig:ch3CA}. Pure biochar exhibited water repellency, whereas the untreated soil was strongly hydrophilic, with a contact angle of only 15.5°. The hydrophobic nature of biochar is commonly attributed to aromatic carbon structures and non-polar aliphatic groups retained or formed during pyrolysis \parencite{pardoMechanismImprovementBiochar2018, maoLinkingHydrophobicityBiochar2019, zhangRolesBiocharsProperties2024}. The contact angle of biochar-amended soil increases by 4° and 8° with 5\% and 10\% biochar addition, respectively.\par
\begin{figure}[htbp]
    \centering
    \includegraphics[width=1.0\linewidth]{../figures/Ch3/ch3CA.png}
    \caption{Contact angle images for (a) biochar, (b) untreated soil, (c) 5\% biochar-amended soil, and (d) 10\% biochar-amended soil}
    \label{fig:ch3CA}
\end{figure}

\subsection{Sample preparation}
To prepare the samples, dry soil and biochar were first weighed and thoroughly mixed, after which de-aired water was added to achieve the OMC. The mixed materials were then sealed in plastic bags and stored for 24 h to ensure moisture equilibration. All the tests were conducted on recompacted specimens. Two compaction states, corresponding to \ac{DC} of 80\% and 95\%, were considered. These compaction levels were selected to represent the relative surface or vegetated soil and densely compacted soil encountered in engineered earthworks, respectively \parencite{gargInfluenceSoilDensity2021, ngEffectsPoreStructure2023, ngPoreStructureEffects2023}. Here, \ac{DC} is defined as the ratio of the dry density to \ac{MDD}. The size of specimens is 80 mm in diameter and 50 mm in height. Specimens were prepared in three layers using the under-compaction method to ensure a uniform density \parencite{laddPreparingTestSpecimens1978, ngEffectsSesquioxideContent2020}. To maintain soil integrity, the surfaces of the first two compacted layers were slightly scarified before adding the soil-biochar mixture \parencite{ngEffectsPoreStructure2023, ngPoreStructureEffects2023}. The as-compacted specimens, kept within the steel sample ring, were submerged in de-aired water and vacuum-saturated for at least 24 h. During the saturation, filter papers were placed on the top and bottom surfaces of the specimens. Two saturation plates (perforated lids) were installed on the tips of the sample ring for protection. Therefore, the volumes of samples can be assumed unchanged after saturation. The samples were weighed to assess their saturated state until a degree of saturation greater than 98\% was achieved. The states of initial and saturated specimens are summarised in Table \ref{tab:ch3_test_programme_states}. Six parallel specimens, marked with an asterisk in Table \ref{tab:ch3_test_programme_states}, were also prepared for the \ac{MIP} and \ac{SEM} tests to investigate the microstructural characteristics of the soil-biochar matrix. These specimens were first saturated. Before testing, carefully trimmed specimens were immersed in liquid nitrogen and then freeze-dried for 72 hours using a Labconco FreeZone freezing dryer to ensure complete drying. Freeze-drying was adopted because it is effective in preserving the original pore structure of the specimens \parencite{ngEffectsPoreStructure2023}.\par
\afterpage{%
\clearpage
\begingroup
\begin{landscape}
\centering
\scriptsize
\setlength{\tabcolsep}{2pt}
\renewcommand{\arraystretch}{1.15}
\setlength{\LTleft}{0pt}
\setlength{\LTright}{0pt}
\newcommand{\dosagecell}[2]{\multirow{##1}{=}{\centering\arraybackslash ##2}}
\begin{longtable}{@{}p{65pt}>{\centering\arraybackslash}p{55pt}>{\centering\arraybackslash}p{56pt}>{\centering\arraybackslash}p{100pt}>{\centering\arraybackslash}p{83pt}>{\centering\arraybackslash}p{56pt}>{\centering\arraybackslash}p{100pt}>{\centering\arraybackslash}p{83pt}@{}}
    \caption{Summary of the specimen details}
    \label{tab:ch3_test_programme_states}\\
    \toprule
    \multirow{2}{=}{Test ID} & \multirow{2}{=}{\makecell[l]{Biochar\\dosage (\%)}} & \multicolumn{3}{c}{After compaction} & \multicolumn{3}{c}{After saturation} \\
    \cmidrule(lr){3-5}\cmidrule(l){6-8}
                                &                                               & \makecell[c]{Void ratio\\($e_0$)} & \makecell[c]{Gravimetric moisture\\content ($w_0$, \%)} & \makecell[c]{Degree of saturation\\($S_{\mathrm{r},0}$, \%)} & \makecell[c]{Void ratio\\($e_{\mathrm{s}}$)} & \makecell[c]{Gravimetric moisture\\content ($w_{\mathrm{s}}$, \%)} & \makecell[c]{Degree of saturation\\($S_{\mathrm{r},\mathrm{s}}$, \%)} \\
    \midrule
    \endfirsthead
    \caption[]{Summary of the specimen details (continued)}\\
    \toprule
    \multirow{2}{=}{Test ID} & \multirow{2}{=}{\makecell[l]{Biochar\\dosage (\%)}} & \multicolumn{3}{c}{After compaction} & \multicolumn{3}{c}{After saturation} \\
    \cmidrule(lr){3-5}\cmidrule(l){6-8}
                                &                                               & \makecell[c]{Void ratio\\($e_0$)} & \makecell[c]{Gravimetric moisture\\content ($w_0$, \%)} & \makecell[c]{Degree of saturation\\($S_{\mathrm{r},0}$, \%)} & \makecell[c]{Void ratio\\($e_{\mathrm{s}}$)} & \makecell[c]{Gravimetric moisture\\content ($w_{\mathrm{s}}$, \%)} & \makecell[c]{Degree of saturation\\($S_{\mathrm{r},\mathrm{s}}$, \%)} \\
    \midrule
    \endhead
    \midrule
    \multicolumn{8}{r}{Continued on next page}\\
    \endfoot
    \bottomrule
    \multicolumn{8}{@{}p{628pt}@{}}{Note: (1) The Test ID indicates the biochar dosage and compaction state. (2) In the Test ID, ``D'' and ``L'' denote dense and loose recompacted specimens, corresponding to DCs of 95\% and 80\%, respectively. (3) The void ratio in the test is assumed to be constant during pre-saturation.}\\
    \endlastfoot
    D-B0*  & \dosagecell{2}{0}  & 0.598 & 13.94 & 63.6 & 0.598 & 21.51 & 98.2 \\
    L-B0*  &                     & 0.895 & 14.65 & 44.7 & 0.895 & 32.78 & 100 \\
    \cmidrule(lr){2-2}
    D-B5*  & \dosagecell{2}{5}  & 0.597 & 15.64 & 66.5 & 0.597 & 23.32 & 99.2 \\
    L-B5*  &                     & 0.893 & 15.3 & 43.5 & 0.893 & 35.16 & 100 \\
    \cmidrule(lr){2-2}
    D-B10* & \dosagecell{2}{10} & 0.611 & 16.47 & 65.2 & 0.611 & 24.92 & 98.7 \\
    L-B10* &                     & 0.906 & 16.33 & 43.6 & 0.906 & 37.44 & 100 \\
\end{longtable}
\end{landscape}
\clearpage
\endgroup
}

\section{Test method}
\subsection{Evaporation test}
A series of evaporation tests was conducted using a combined tensiometer and dew-point hygrometry method \parencite{satyanagaMeasurementSoilwaterCharacteristic2019, changWaterTransferGranite2026}. As shown in Figure \ref{fig:ch3WP4CHyprop2}a, a water evaporation system (METER Group) consists of tensiometer sensors (resolution 0.001 kPa), a specimen ring, an electronic balance (resolution: 0.01 g), and a computer for data acquisition. The bottom of the specimen is sealed to maintain a zero-flux boundary condition while the water is only allowed to evaporate from its top surface. The mass of evaporating water is continuously measured and recorded by the electronic balance. Two vertically aligned tensiometer shafts are installed within the specimen to continuously monitor matric suction, with their ceramic tips positioned 12.5 mm above and below the specimen centre, respectively. The matric suction of the specimen was calculated as the mean value measured by the two tensiometers. The maximum matric suction that can be reliably measured by the water evaporation system is approximately 100 kPa. To extend the measurement of suction range, a WP4C device based on the chilled-mirror dew-point technique is employed \parencite{shokranaMeasurementSoilWater2020, changWaterTransferGranite2026}. It should be noted that the WP4C measures total suction, which comprises both matric and osmotic suction components. However, the osmotic suction is generally negligible for the soil tested in this study, and previous studies have demonstrated good agreement between WP4C and tensiometer measurements \parencite{aziziCouplingCyclicWater2023, dongHydrocharEffectsMicrostructure2025a}.\par
The unsaturated \ac{HC} is calculated with the simplified evaporation model \parencite{schindlerSimplifyingEvaporationMethod2006, petersSimplifiedEvaporationMethod2008}. Assuming the matric suction (i.e., negative pore-water pressure) varies continuously and linearly along the vertical direction of the specimen, a constant \ac{HC} can be determined at a given time during the transient process. According to Darcy’s law, \ac{HC} is given by:
\begin{equation}
    k_{un}=-\frac{\Delta V}{2A(t_{i+1}-t_{i})} \cdot \frac{1}{\nabla s}
\end{equation}
where $t_{i}$ is the i-th elapsed time, $\Delta V$ is the volume of evaporating water over the time interval, $A$ is the cross-sectional area of the specimen, and $s$ is matric suction. The average suction gradient, $\nabla s$4, over the time interval ($t_{i}$, $t_{i+1}$) can be expressed as follows:
\begin{equation}
    \nabla s = \frac{1}{2 \gamma_{\mathrm{w}} \cdot \Delta h}(\Delta s_{i+1} + \Delta s_{i}) - 1
\end{equation}
where $\gamma_{\mathrm{w}}$ is the unit weight of water, $\Delta h$ is the vertical spacing between the two ceramic tips, taken as 25 mm, and $\Delta s_{i}$ and $\Delta s_{i+1}$ are the differences in matric suctions measured by the upper and lower tensiometers at $t_{i}$ and $t_{i+1}$, respectively.\par
The \ac{SWRC} curves were determined by combining continuous measurements from the water evaporation system and discrete measurements from WP4C \parencite{satyanagaMeasurementSoilwaterCharacteristic2019, singhEstimatingSoilWater2024}. Before the evaporation test, a set of syringes was used to evacuate trapped air from the ceramic tips and sensor ports (Figure \ref{fig:ch3WP4CHyprop2}b) until no visible bubbles were observed. This aims to make them fully saturated before installation. During the test, the specimen was allowed to evaporate naturally under laboratory conditions at approximately 22 °C and 60\% relative humidity. When the measured matric suction approached the cavitation limit of the ceramic tips at 100 kPa, the test was transferred to the WP4C device. The specimen was then removed from the water evaporation system and carefully cut into blocks of approximately 2 cm \parencite{changWaterTransferGranite2026}. These sample blocks were then subjected to continued evaporation under the same laboratory conditions. Three selected blocks, taken from the bottom, middle, and upper layers, respectively, were tested using WP4C to determine their water potential, after which their water contents were measured immediately \parencite{dongMicrostructuralInvestigationUnsaturated2024}. The mean water potential and water content of the three blocks were used to determine \ac{SWRC}.\par
\begin{figure}[htbp]
    \centering
    \includegraphics[width=1.0\linewidth]{../figures/Ch3/ch3WP4C&Hyprop2.png}
    \caption{Experimental systems: (a) water evaporation testing system and dew-point hygrometry, and (b) saturation of the testing apparatus}
    \label{fig:ch3WP4CHyprop2}
\end{figure}

\subsection{Microstructure and pore structure tests}
\Ac{SEM} test was performed using a Philips XL30 ESEM in secondary electron mode to examine the microstructure and interfacial interactions between soil and biochar. Before testing, samples were gold-coated to enhance conductivity. The \ac{PSD} was determined by \ac{MIP} test using an AutoPore V9600 porosimeter (Micromeritics Inc., USA). The cumulative volume and mercury intrusion pressure were recorded at each intrusion step. By idealising the mercury-accessible pore entrance as cylindrical capillaries, the equivalent pore-throat radius can be determined using the Washburn equation \parencite{washburnDynamicsCapillaryFlow1921}.

\section{Results and analysis}
\subsection{Pore size distribution}
Figure \ref{fig:ch3PoreSize} presents the logarithmic differential PSD ($\Delta V/\Delta \mathrm{log}_{10}D$) and cumulative mercury intrusion curves obtained from the MIP test. In the current work, the soil pore system is divided into macropores (> 75 $\mu$m), mesopores (30–75 $\mu$m), micropores (5–30 $\mu$m), ultra-micropores (0.1–5 $\mu$m), and crypto-pores (< 0.1 $\mu$m) ranges adopted from Dong et. al \parencite*{dongMicrostructuralInvestigationUnsaturated2024}. Different pore ranges are represented with different background colours in Figure \ref{fig:ch3PoreSize} for clarity. As depicted in Figure \ref{fig:ch3PoreSize}a and c, the tested specimens exhibit bimodal PSD curves, with the first pore peak ($D_{\mathrm{m}}$) in the ultra-micropore range, and a second pore peak ($D_{\mathrm{s}}$) in the micro-to-macropore ranges. A pronounced valley is observed between the two peaks, and the corresponding pore diameter is defined as a delimiting pore size ($D_{\mathrm{L}}$). A summary of $D_{\mathrm{m}}$, $D_{\mathrm{s}}$, $D_{\mathrm{L}}$, and mercury-detected void ratio is listed in Table \ref{tab:ch3_pore_structure_parameters}.
For densely compacted specimens, biochar has little influence on the \ac{PSD} within the cryptopore range. In the ultra-micropore range, with biochar addition, $D_{\mathrm{m}}$ increases from 0.43 $\mu$m to 2.05 $\mu$m and 2.07 $\mu$m for D-B5 and D-B10, respectively. The pore volume within the 0.5-3 $\mu$m range shows a significant increase with biochar addition, consistent with the micropore expansion effect reported by Dong et al. \parencite*{dongMicrostructuralInvestigationUnsaturated2024}. For D-B5, a weak residual peak remains at approximately 0.35 $\mu$m, close to the $D_{\mathrm{m}}$ of D-B0. But this feature disappears in D-B10, where the PSD curve increases continuously towards its first peak. Despite doubling the biochar content, $D_{\mathrm{m}}$ remained unchanged in both position and intensity for D-B5 and D-B10, together with almost overlapping descending branches from $D_{\mathrm{m}}$ to $D_{\mathrm{L}}$. In addition, the values of $D_{\mathrm{L}}$ for all three densely compacted specimens are the same at around 5 $\mu$m. These investigations suggest that $D_{\mathrm{m}}$ is unlikely to be governed by the intrinsic pores of the biochar, as a marked change would otherwise be expected when the biochar dosage is doubled. Instead, it is associated with structural pores formed between biochar and soil particles, as further discussed based on the \ac{SEM} results in the following section. Micropores and mesopores are most strongly affected by biochar. In these ranges, $\Delta V/\Delta \mathrm{log}_{10}D$ of D-B5 is approximately one-third of that of D-B0, while with increasing biochar dosage to 10\%, a partial recovery of pore value is observed. In contrast, \ac{MIP}-detected macropores only changed slightly with biochar. Figure \ref{fig:ch3porecategory} presents the proportions of different pore categories detected by the \ac{MIP} method. At 95\% \ac{DC}, the addition of biochar introduces more ultra-micropores and fewer micropores. Similar behaviour has also been observed in previous studies \parencite{yangPyrolysisTemperatureAffects2021, ngEffectsSoilPlant2022, yiEffectsFirwoodBiochar2024}. But the proportions of other pore categories only change slightly.\par
Compared with 95\% \ac{DC} specimens, the \ac{PSD} curves of 80\% \ac{DC} specimens exhibit a more pronounced bimodal pattern. For L-B0, $D_{\mathrm{m}}$ is centred at 0.82 $\mu$m. With biochar, $D_{\mathrm{m}}$ increases to 2.15 $\mu$m and 2.13 $\mu$m for L-B5 and L-B10, respectively. Similar to densely compacted specimens, biochar increases the pore volume between $D_{\mathrm{m}}$ of L-B0 and 5 $\mu$m. However, the effect of biochar on larger pore sizes is different. Biochar enhances the volume of 5-30 $\mu$m pores but decreases that of mesopores and 75-156 $\mu$m macropores. Only those macropores larger than 156 $\mu$m increase with larger biochar dosage. The value of $D_{\mathrm{s}}$ decreases with more biochar. From Figure \ref{fig:ch3porecategory}, only micropores show a slight increase with biochar dosage, while the proportions of other pore categories remain relatively stable.\par
Figures \ref{fig:ch3PoreSize}b and d present the cumulative mercury intrusion curves. Reducing the \ac{DC} results in a higher cumulative intrusion volume. From Table Table \ref{tab:ch3_pore_structure_parameters}, the \ac{MIP}-intruded void ratio is 3-12\% lower than the total void ratio calculated from the specimen mass-volume relationships. This difference represents non-intruded porosity, which may comprise pores outside the measurable size range of the porosimeter and closed or isolated pores that are not externally accessible to mercury \parencite{romeroMicrostructureInvestigationUnsaturated2008, elmountassirExperimentalStudyCompaction2014}. The densely compacted specimens show a greater proportion of \ac{MIP}-inaccessible pores, likely due to compaction-induced narrowing of pore entrances, producing stronger pore-body–throat contrasts and ink-bottle effects \parencite{delageStudyStructureSensitive1984, yuanXRayMicrotomographyMercury2022}.\par
\begin{figure}[htbp]
    \centering
    \includegraphics[width=1.0\linewidth]{../figures/Ch3/ch3PoreSize.png}
    \caption{Measured pore size distribution from MIP test of densely compacted specimens (a and b), and loosely compacted specimens (c and d)}
    \label{fig:ch3PoreSize}
\end{figure}
\begin{figure}[htbp]
    \centering
    \includegraphics[width=1.0\linewidth]{../figures/Ch3/ch3porecategory.png}
    \caption{Percentage of different pore categories in the tested specimens}
    \label{fig:ch3porecategory}
\end{figure}
\begin{table}[htbp]
    \centering
    \caption{Pore structure parameters of representative specimens}
    \label{tab:ch3_pore_structure_parameters}
    \begingroup
    \renewcommand{\arraystretch}{1.15}
    \setlength{\tabcolsep}{2pt}
    \begin{tabular*}{\textwidth}{@{\extracolsep{\fill}}>{\centering\arraybackslash}m{0.12\textwidth}>{\centering\arraybackslash}m{0.11\textwidth}>{\centering\arraybackslash}m{0.15\textwidth}>{\centering\arraybackslash}m{0.18\textwidth}>{\centering\arraybackslash}m{0.10\textwidth}>{\centering\arraybackslash}m{0.10\textwidth}>{\centering\arraybackslash}m{0.10\textwidth}@{}}
        \toprule
        \makecell[c]{Specimen\\ID} & \makecell[c]{Void\\ratio} & \makecell[c]{Intruded\\void ratio} & \makecell[c]{Non-detected\\void ratio} & \makecell[c]{$D_{\mathrm{m}}$\\($\mu$m)} & \makecell[c]{$D_{\mathrm{L}}$\\($\mu$m)} & \makecell[c]{$D_{\mathrm{s}}$\\($\mu$m)} \\
        \midrule
        D-B0  & 0.598 & 0.531 & 0.067 & 0.43 & 5.33 & 18.87 \\ 
        L-B0  & 0.895 & 0.822 & 0.073 & 0.82  & 6.52 & 60.05 \\ 
        D-B5  & 0.597  & 0.555 & 0.042 & 2.05  & 5.17 & 44.91 \\ 
        L-B5  & 0.893 & 0.857 & 0.036 & 2.15  & 7.17 & 29.84 \\ 
        D-B10 & 0.611 & 0.552 & 0.059 & 2.07  & 4.82 & 30.62 \\ 
        L-B10 & 0.906 & 0.88  & 0.026 & 2.13  & 6.52 & 23.27 \\ 
        \bottomrule
    \end{tabular*}
    \endgroup
\end{table}

\subsection{Microstructural characteristics}
Figure \ref{fig:ch3SEM} presents the \ac{SEM} images of the tested specimens. In specimen D-B0, the soil particles form aggregates, which generate pronounced macro- and mesopores between adjacent aggregates, i.e., inter-aggregate pores \parencite{muHydromechanicalBehaviorUnsaturated2023}. These pores are interconnected through smaller pore throats, forming preferential pathways for water flow, as illustrated in Figure \ref{fig:ch3SEM}a. These pore throats can induce an ink-bottle structure \parencite{phamStudyHysteresisModels2005, ngWaterRetentionVolumetric2016}, explaining smaller MIP-detected void ratios in Table \ref{tab:ch3_pore_structure_parameters}. A higher-magnification image focusing on a selected aggregate is presented in Figure \ref{fig:ch3SEM}b. The image shows that aggregates mainly contain intra-aggregate pores, which have characteristics of sizes less than 5 $\mu$m. Therefore, it can be interpreted that intra-aggregate pores are associated with $D_{\mathrm{m}}$, while inter-aggregate pores correspond to $D_{\mathrm{s}}$. At a lower \ac{DC} (Figure \ref{fig:ch3SEM}c), the soil particles still form aggregates, but the inter-aggregate pores are more noticeable. This investigation is consistent with Figure \ref{fig:ch3porecategory}, which demonstrates that compaction mainly reduces meso- and macropores.\par
With biochar addition, the sand, silt and biochar particles in L-B5 form an interlocked structure with predominantly face-to-face contacts (Figure \ref{fig:ch3SEM}e), which can enhance aggregate stability \parencite{baiamonteEffectBiocharPhysical2019, muHydromechanicalBehaviorUnsaturated2023}. The sand particles are sub-rounded and in close contact with the biochar particles, while the intra-aggregate pores are partially filled by silt particles and biochar fragments. This explains the increase in certain ultra-micropore volumes following biochar addition observed in Figures \ref{fig:ch3PoreSize}a and c. At higher magnification (Figure \ref{fig:ch3SEM}f), mesopores and irregular macropores are observed at the biochar-soil interfaces. By comparison, D-B5 in Figure \ref{fig:ch3SEM}d shows fewer mesopores and macropores than those of L-B5, producing a more complex and heterogeneous pore network with tortuous patterns. Comparison among Figures \ref{fig:ch3SEM}c, f and i shows that the heterogeneous distribution of biochar particles generates additional inter-aggregate macropores, while fine biochar fragments fill some of the mesopores. These pores can provide more effective water flow paths \parencite{blanco-canquiBiocharSoilPhysical2017}. This explains that as biochar dosage increases, $D_{\mathrm{s}}$ decreases and a greater macropore proportion is observed in Figure \ref{fig:ch3PoreSize}c. Figures \ref{fig:ch3SEM}g and h indicate the contact between the clay and biochar. On the one hand, clay particles invade biochar intrinsic pores, partially occluding the internal pore network \parencite{jingInteractionsBiocharClay2022}. On the other hand, fractured intrinsic pores can provide attachment sites for clay particles due to electrostatic attraction \parencite{yaoCharacterizationEnvironmentalApplications2014}. Owing to the hydrophilic nature of clay (Figure \ref{fig:ch3CA}), clay particles attached to the biochar surface may increase its affinity for water.\par
\begin{figure}[htbp]
    \centering
    \includegraphics[width=1.0\linewidth]{../figures/Ch3/ch3SEM.png}
    \caption{\Ac{SEM} images of the tested materials: (a) D-B0 at ×500, (b) magnified view at×2000, (c) L-B0 at ×500, (d) D-B5 at ×500, (e) L-B5 at ×100, (f) magnified view at ×500, (g) L-B5 at ×2000, (h) L-B5 at ×2000, and (i) L-B10 at ×500}
    \label{fig:ch3SEM}
\end{figure}

\subsection{Soil water retention characteristic}
The experimental data of SWRC were fitted using the van Genuchten (VG) model \parencite{vangenuchtenClosedformEquationPredicting1980}, which can be expressed as:
\begin{equation}
    \theta(s) = \theta_{\mathrm{r}} + \frac{\theta_{\mathrm{s}} - \theta_{\mathrm{r}}}{[1 + (\alpha \cdot S)^n]^m}
\end{equation}
where $\theta_{\mathrm{s}}$ is the saturated volumetric water content, $\theta_{\mathrm{r}}$ is the residual volumetric water content, and $\alpha$, $m$, and $n$ are the fitting parameters. The volumetric water content is defined as the ratio of water volume to the total volume of the sample. Figure \ref{fig:ch3SWRC} presents the SWRC results. The Casagrande method was used to determine the \ac{AEV} \parencite{vanapalliInfluenceSoilStructure1999}. The fitting parameters and \acp{AEV} are summarised in Table \ref{tab:ch3_vg_fitting_parameters}.\par
For both densely and loosely compacted specimens, the addition of biochar causes only minor changes in the \ac{SWRC} in the near-saturated range, particularly at matric suctions below 10 kPa. With increasing matric suction, the biochar-amended specimens exhibit enhanced water retention capacity. A comparison between Figures \ref{fig:ch3SWRC}a and b further shows that biochar has a more pronounced effect on the \ac{SWRC} under denser compaction than under loose compaction in a relatively lower matric suction range. This can be interpreted from the \ac{SEM} results. The densely compacted soil contains fewer inter-aggregate mesopores and macropores, while the loosely compacted soil contains a relatively developed inter-aggregate pore network. In addition, increasing the biochar content leads to a more pronounced modification of pore structure in densely compacted samples, making the corresponding \ac{SWRC} more sensitive to biochar dosage.\par
At 95\% \ac{DC}, the \ac{AEV} increases from 4.986 kPa for the untreated soil to 6.152 kPa for the 5\% biochar-amended sample but subsequently decreases to 5.017 kPa at a biochar content of 10\%. By contrast, at 80\% \ac{DC}, the \ac{AEV} increases progressively from 3.017 kPa for L-B0 to 3.894 kPa for L-B5 and 4.285 kPa for L-B10. According to the Young-Laplace equation, the equivalent pore diameter corresponding to the \ac{AEV} can be calculated as:
\begin{equation}
    D_{\mathrm{AEV}} = \frac{4 \gamma \cdot \cos \varphi}{s_{\mathrm{AEV}}}
\end{equation}
where $\varphi$ is the contact angle, the values of which can be taken from Figure \ref{fig:ch3CA}. As listed in Table \ref{tab:ch3_vg_fitting_parameters}, the values of $D_\mathrm{AEV}$ indicate that the measured \acp{AEV} are primarily associated with pore throats of approximately 40-100 $\mu$m. The corresponding \ac{PSD} therefore provides a basis for interpreting the changes in \ac{AEV}. For the specimens compacted to 95\% \ac{DC}, the \ac{PSD} intensity below approximately 50 $\mu$m follows the order D-B0 > D-B10 > D-B5, consistent with the \ac{AEV} increasing in the reverse order (D-B0 < D-B10 < D-B5). The lower \ac{AEV} at 10\% than at 5\% biochar therefore reflects a non-monotonic restructuring of the controlling pore throats under denser compaction. In contrast, at 80\% \ac{DC}, biochar progressively reduces the abundance of pore throats within the relevant size range, restricting air entry and resulting in a monotonic increase in \ac{AEV} with biochar content. As indicated by the \ac{SEM} observations, the specimens at lower compaction contain more meso- and macropores. Biochar has a relatively limited influence on these pre-existing inter-aggregate pores, resulting in only minor differences among the \ac{SWRC} curves of the biochar-amended specimens in Figure \ref{fig:ch3SWRC}b.\par
In the VG model, the fitting parameter $n$ controls the dehydration rate of the \ac{SWRC} curve over the desaturation range. The value of $n$ is reported to reflect the characteristics of the soil \ac{PSD} \parencite{ngWaterRetentionVolumetric2016}. As summarised in Table \ref{tab:ch3_vg_fitting_parameters}, the values of $n$ are consistently higher at 80\% DC than at 95\% \ac{DC}. This difference is attributed to the greater proportion of meso- and macropores in the more loosely compacted specimens. It also noted that the value of $n$ decreases with biochar. This can be attributed to less pore volume induced by biochar amendment \parencite{hussainInvestigatingUnsaturatedHydraulic2021, hussainInvestigatingBiocharamendedSoil2022}. However, the effect of biochar on $n$ is more complex. At 80\% \ac{DC}, $n$ decreases monotonically with increasing biochar content, whereas at 95\% \ac{DC} it initially decreases and then increases. This phenomenon is the same as the change in \ac{AEV}. As shown by the \ac{PSD} results, the \ac{PSD} above $D_\mathrm{L}$ changes monotonically with increasing biochar content at 80\% \ac{DC}, but non-monotonically at 95\% \ac{DC}. This further indicates that biochar induces compaction-dependent pore redistribution rather than a uniform shift of the entire pore system towards finer pores. This effect is discussed in Section \ref{sec:ch3discussion}.\par
\begin{figure}[htbp]
    \centering
    \includegraphics[width=1.0\linewidth]{../figures/Ch3/ch3HYPROPWP4CSWRC.png}
    \caption{\Ac{SWRC} in the evaporation tests: (a) densely compacted specimens, and (b) loosely compacted specimens}
    \label{fig:ch3SWRC}
\end{figure}

\begin{table}[htbp]
    \centering
    \caption{Fitting parameters of the VG model and \ac{AEV}}
    \label{tab:ch3_vg_fitting_parameters}
    \begingroup
    \scriptsize
    \renewcommand{\arraystretch}{1.15}
    \setlength{\tabcolsep}{2pt}
    \begin{tabular*}{\textwidth}{@{\extracolsep{\fill}}ccccccccc@{}}
        \toprule
        \makecell[c]{Specimen\\ID} & $\theta_{\mathrm{r}}$ & $\theta_{\mathrm{s}}$ & \makecell[c]{$\alpha$\\(kPa$^{-1}$)} & $n$ & $m$ & $R^2$ & \makecell[c]{AEV\\(kPa)} & \makecell[c]{Equivalent pore-throat\\diameter ($\mu$m)} \\
        \midrule
        D-B0  & 0.0853 & 0.374 & 0.105 & 1.128 & 0.113 & 0.994 & 4.986 & 57.34 \\
        L-B0  & 0.0852 & 0.445 & 0.144 & 1.24  & 0.194 & 0.998 & 3.017 & 93 \\
        D-B5  & 0.116  & 0.368 & 0.085 & 1.125 & 0.111 & 0.995 & 6.152 & 44.62 \\
        L-B5  & 0.0851 & 0.445 & 0.136 & 1.175 & 0.149 & 0.997 & 3.894 & 70.47 \\
        D-B10 & 0.085  & 0.369 & 0.104 & 1.077 & 0.072 & 0.992 & 5.015 & 53.25 \\
        L-B10 & 0.0852 & 0.447 & 0.121 & 1.162 & 0.139 & 0.995 & 4.285 & 62.35 \\
        \bottomrule
    \end{tabular*}
    \endgroup
\end{table}

\subsection{Unsaturated hydraulic conductivity}
Figure \ref{fig:ch3HC} presents the relationship between the unsaturated \ac{HC} and matric suction. The variation of unsaturated \ac{HC} ($k_\mathrm{un}$) with the matric suction ($s$) is fitted using the van Genuchten-Mualem (VGM) model \parencite{mualemNewModelPredicting1976, vangenuchtenClosedformEquationPredicting1980}, which can be expressed as:
\begin{equation}
    k_\mathrm{un}(s) = k_0 \cdot S_\mathrm{e}^{l} \cdot [1 - (1 - S_\mathrm{e}^{1/m})^m]^2
\end{equation}
where $k_0$ is the saturated hydraulic conductivity, $l$ is a pore-connectivity parameter and can be taken as an empirical value of 0.5, and $S_\mathrm{e}$ is the effective degree of saturation, which is defined as:
\begin{equation}
    S_\mathrm{e}(s) = \frac{\theta(s) - \theta_\mathrm{r}}{\theta_\mathrm{s} - \theta_\mathrm{r}} = [1 + (\alpha \cdot s)^n]^{-m}
\end{equation}
The values of $m$, $n$, $\alpha$, $\theta_\mathrm{s}$, and $\theta_\mathrm{r}$ are adopted from the VG model, while only $k_0$ is fitted. To validate the fitting results, the values of fitted $k_0$ are compared with the experimental results, which are tested via the falling-head method following BS EN ISO 17892-11. The fitted and tested $k_0$ are denoted by $k_{0,\mathrm{f}}$ and $k_{0,\mathrm{t}}$, respectively. The results are summarised in Table \ref{{tab:ch3_vgm_fitting_parameters}}. It can be observed that the values of $k_{0,\mathrm{f}}$ and $k_{0,\mathrm{t}}$ are in good agreement, demonstrating the reliability of the fitted parameters.\par
At matric suctions below 3 kPa, the specimens compacted to 80\% \ac{DC} exhibit higher $k_\mathrm{un}$ than those compacted to 95\% \ac{DC} owing to their larger fractions of macro- and mesopores. Among the 80\% \ac{DC} specimens, L-B0 has the largest secondary pore peak and the highest combined macro- and mesopore fraction (approximately 30\%), corresponding to the highest $k_\mathrm{un}$. However, its unsaturated \ac{HC} decreases most rapidly once the matric suction exceeds the \ac{AEV}. Accordingly, the values of $k_\mathrm{un}$ rank in the order L-B0 > L-B5 > L-B10. A similar correspondence between pore structure and unsaturated \ac{HC} is observed at 95\% \ac{DC}. The combined macro- to micropore fraction follows the order D-B0 > D-B10 > D-B5, consistent with the order of $k_\mathrm{un}$. These results indicate that, in the low-suction range, unsaturated \ac{HC} is governed mainly by the abundance and connectivity of the larger drainage pores. For all the tested specimens, $k_\mathrm{un}$ decreases rapidly after the matric suction exceeds their \acp{AEV} (corresponding equivalent pore throat around 75 $\mu$m). This is because air first invades the larger pores, reducing the effective cross-sectional area available for water flow and increasing the tortuosity of the remaining flow paths \parencite{vangenuchtenClosedformEquationPredicting1980, fredlundPredictingPermeabilityFunction1994}. Driven by capillary pressure, water tends to form thin films coating soil and biochar particles or to be retained within smaller pores \parencite{anEffectsSoilCharacteristics2018}. Therefore, continued drying reduces both the thickness and hydraulic conductance of these water films, thereby increasing viscous resistance and decreasing $k_\mathrm{un}$ \parencite{tullerHydraulicConductivityVariably2001}.\par
As water drains from the micro-to-macropores (< 5 $\mu$m), $k_\mathrm{un}$ becomes increasingly governed by the finer pore network. Under both \acp{DC}, the dominant pore diameters of the biochar-amended specimens are very similar (2.05 and 2.07 $\mu$m for the 95\% \ac{DC} specimens and 2.15 and 2.13 $\mu$m for the 80\% \ac{DC} specimens), which is consistent with the near convergence of B5 and B10 $k_\mathrm{un}$ curves beyond around 60 kPa. The slightly higher $k_\mathrm{un}$ of B10 at high suctions can be explained by its more water flow paths observed in the \ac{SEM} test (Figure \ref{fig:ch3SEM}i).\par
\begin{figure}[htbp]
    \centering
    \includegraphics[width=1.0\linewidth]{../figures/Ch3/ch3HC.png}
    \caption{Unsaturated \ac{HC} of the specimens at (a) 95\% \ac{DC} and (b) 80\% \ac{DC}}
    \label{fig:ch3HC}
\end{figure}

\begin{table}[htbp]
    \centering
    \caption{Fitting parameters of the VGM model}
    \label{tab:ch3_vgm_fitting_parameters}
    \begingroup
    \renewcommand{\arraystretch}{1.15}
    \setlength{\tabcolsep}{2pt}
    \begin{tabular*}{\textwidth}{@{\extracolsep{\fill}}ccccc@{}}
        \toprule
        \makecell[c]{Specimen\\ID} & \makecell[c]{$k_{0,\mathrm{f}}$\\($\times 10^{-7}$ m/s)} & $R^2$ & \makecell[c]{$k_{0,\mathrm{t}}$\\($\times 10^{-7}$ m/s)} & $k_{0,\mathrm{t}}/k_{0,\mathrm{f}}$ \\
        \midrule
        D-B0  & 2.036 & 0.945 & 2.402 & 1.18 \\
        L-B0  & 2.617 & 0.995 & 3.873 & 1.48 \\
        D-B5  & 1.16  & 0.912 & 1.53  & 1.32 \\
        L-B5  & 2.931 & 0.957 & 4.133 & 1.41 \\
        D-B10 & 2.931 & 0.95  & 4.483 & 1.53 \\
        L-B10 & 2.405 & 0.974 & 3.271 & 1.36 \\
        \bottomrule
    \end{tabular*}
    \endgroup
\end{table}

\section{Discussion}
\label{sec:ch3discussion}
\subsection{Fractal dimension analysis}
\label{sec:fractaldimension}
For a three-dimensional pore system, the fractal dimension of the pore surface ($d_\mathrm{s}$) can characterise the self-similarity and complexity of the pore system \parencite{perfectApplicationsFractalsSoil1995, hanRelationshipFractalFeature2022}. In this work, ds is determined using a thermodynamic fractal model, which is suitable for describing the pore structure of geotechnical materials \parencite{shiStrengthMicroscopicPore2024}. By assuming the pore system as an equivalent bundle of ideal capillary tubes, the relationship between the external pressure and the incremental volume of intruded mercury can be expressed as \parencite{shiStrengthMicroscopicPore2024}:
\begin{equation}
    \log_{10} \left[ \left( \sum_{i = 1}^{n} p_{\mathrm{m},i} \Delta V_i \right) \middle/ r_{n}^2 \right] = d_{\mathrm{s}} \log_{10} \left( V_{n}^{1/3} \middle/ r_{n} \right) + R
\end{equation}
where $p_{\mathrm{m},i}m$ is the applied pressure at the i-th intrusion step, $\Delta V_i$ represents the incremental volume of mercury at the i-th intrusion step, $V_n$ denotes the accumulative volume of mercury at the n-th intrusion step, $r_n$ is the equivalent radius of the pore, $d_s$ and $R$ can be determined as the slope and intercept in the linear regression, respectively, as presented in Figure \ref{fig:ch3Fractalline}. The results of fractal dimension are summarised in Figure \ref{fig:ch3Fractalvalue}.\par
At all biochar dosages, $d_s$ is higher at 95\% \ac{DC} than at 80\% \ac{DC}, indicating a rougher and geometrically more complex pore-solid interface. A similar density dependence was reported by Sun et al. \parencite*{sunFractalCharacteristicsPore2020}, which observed a slight increase in the thermodynamic fractal dimension with increasing dry density and showed that compaction primarily affected the macropore domain, but intra-aggregate pores remined nearly unchanged. At 95\% \ac{DC}, ds increases at 5\% biochar and then decreases slightly at 10\%, consistent with the variation in the secondary \ac{PSD} peak at $D_\mathrm{s}$ (Figure \ref{fig:ch3PoreSize}a). The maximum ds at 5\% indicates the greatest pore-structure heterogeneity, associated with the coexistence of residual inter-aggregate pores, soil-biochar interfacial voids and finer pores. Similarly, three-dimensional micro-CT observations in a previous study also show that biochar amendment can increase pore irregularity and structural complexity \parencite{sunThreedimensionalFractalCharacteristics2021}. In contrast, at 80\% \ac{DC}, the $D_\mathrm{m}$ and $D_\mathrm{s}$ peaks progressively converge with increasing biochar content, indicating reduced differentiation between the two pore domains and accounting for the monotonic decrease in $d_s$.\par
\begin{figure}[htbp]
    \centering
    \includegraphics[width=1.0\linewidth]{../figures/Ch3/ch3FractalLine.png}
    \caption{Regression analysis of fractal dimensions of the tested specimens: (a) 95\% \ac{DC}, and (b) 80\% \ac{DC}}
    \label{fig:ch3Fractalline}
\end{figure}

\begin{figure}[htbp]
    \centering
    \includegraphics[width=0.56\linewidth]{../figures/Ch3/ch3fractalDimensionValue.png}
    \caption{Fractal dimension values of the tested specimens}
    \label{fig:ch3Fractalvalue}
\end{figure}

\subsection{Pore-domain-specific hydraulic response}
Based on the above investigations, the pore classes are regrouped into three functional integration domains: a fine retention domain (<5 $\nu$m) dominated by intra-aggregate pores, a drainage domain (5-75 $\nu$m) with larger intra-aggregate pores and smaller inter-aggregate throats, and a large inter-aggregate/interfacial domain (>75 $\nu$m). To compare the soil water retention capacity over a specified matric suction range ($s_1$, $s_2$), an integral mean water content (\ac{IMWC}) is first defined as \parencite{romanoRootZoneWaterStorageCapacity2025}:
\begin{equation}
    \mathrm{IMWC} = \frac{1}{\log_{10} \left(s_2 \middle/ s_1 \right)} \int_{s_1}^{s_2} \,d\log_{10} \left( s \right)
\end{equation}
For the biochar-amended samples, the change in \ac{IMWC} related to the untreated soil with the identical \ac{DC} can be estimated by:
\begin{equation}
    I_\mathrm{WR} = \mathrm{IMWC}_\mathrm{b} - \mathrm{IMWC}_0
\end{equation}
where $\mathrm{IMWC}_\mathrm{b}$ and $\mathrm{IMWC}_0$ are calculated from the biochar-amended and untreated specimens, respectively. Similarly, an integral mean hydraulic conductivity (\ac{IMHC}) is defined to quantify the average hydraulic conductivity over ($s_1$, $s_2$):
\begin{equation}
    \mathrm{IMHC} = \frac{1}{\log_{10} \left( s_2 \middle/ s_1 \right)} \int_{}^{}  \,d\log_{10} \left( s\right)
\end{equation}
The biochar-induced HC penalty is given by:
\begin{equation}
    I_\mathrm{k} = \mathrm{IMWC}_\mathrm{0} - \mathrm{IMWC}_\mathrm{b}
\end{equation}
where $IMHC_\mathrm{b}$ and $IMHC_0$ are calculated from biochar-treated and untreated samples, respectively. It should be noted that the \acp{SWRC} were fitted using the VG model over a wider suction range, whereas the \acp{HC} were fitted using the VGM model only within the range of 0.01-100 kPa, owing to the measurement limit of the tensiometers in the water evaporation system. Different suction ranges are illustrated in Figure \ref{fig:ch3IWRIk}a. Therefore, $I_\mathrm{WR}$ and $I_\mathrm{k}$ were initially calculated over suction ranges of 0.01-10,000 kPa and 0.01-100 kPa (Figures \ref{fig:ch3IWRIk}b and c), respectively. To further examine their relationship, both indices were evaluated over their common suction range of 0.01-100 kPa in Figure \ref{fig:ch3IWRIk}d.\par
Under denser compaction, biochar reduces the contribution of the >75 $\mu$m pore domain to water retention ($I_\mathrm{WR}$ < 0). As discussed in Section \ref{sec:fractaldimension}, the addition of biochar at 95\% \ac{DC} promotes the formation of interfacial pores. These pores provide additional drainage pathways during drying, thereby reducing $I_\mathrm{WR}$. At 80\% \ac{DC}, however, the reduction in the \ac{MIP}-detected >75 $\mu$m pore volume is accompanied by only small increases in $I_\mathrm{WR}$ of 0.0018 and 0.0039 (Figure \ref{fig:ch3IWRIk}b). All amended specimens also show positive values of $I_\mathrm{k}$ (0.040–0.234). The drainage domain shows the clearest compaction dependence in water storage, with $I_\mathrm{WR}$ lower in dense specimens (0.0023-0.0073) than in loose specimens (0.0207-0.0291). Nevertheless, $I_\mathrm{k}$ remains positive for all treatments (0.037-0.107), corresponding to \ac{HC} reductions of 8-22\%. Under 95\% \ac{DC}, the response is non-monotonic, with a greater conductivity penalty for D-B5 than for D-B10. This divergence between storage and conductivity indicates that biochar-induced pore redistribution does not translate directly into hydraulic transmission. Instead, the hydraulic response depends on how the initial fabric constrains the connectivity, throat size and tortuosity of the restructured drainage network. The retention domain provides the largest increase in water storage, particularly under loose compaction. From Figure \ref{fig:ch3IWRIk}c, the unsaturated \ac{HC} in retention domain is significantly enhanced, suggesting that fine pores and water films can preserve liquid continuity after water drainage from larger pores.\par
Figure \ref{fig:ch3IWRIk}d illustrates the integrated effects of biochar content and \ac{DC} on soil hydraulic behaviours. The hydraulic response to biochar exhibits a clear compaction dependence and an inherent trade-off between water retention capacity and \ac{HC}, particularly within the drainage domain. Under denser compaction, the limited macropore storage is accompanied by reduced drainage and lower conductivity ($I_\mathrm{WR}$ < 0 and $I_\mathrm{k}$ > 0), indicating that the remaining drainage pathways are relatively restricted. With less \ac{DC}, biochar increases water storage within the larger-pore domain while still reducing \ac{HC} ($I_\mathrm{WR}$ > 0 and $I_\mathrm{k}$ > 0), suggesting that additional storage does not necessarily translate into more efficient water transmission. In the retention domain, however, biochar can simultaneously increase water storage and facilitate residual water flow ($I_\mathrm{WR}$ >0 and $I_\mathrm{k}$ <0), probably by preserving liquid continuity within fine pores and water films after the larger pores have drained. Nevertheless, no enhancement in HC at the expense of water-retention capacity is observed across the tested suction range and biochar dosages.\par
\begin{figure}[htbp]
    \centering
    \includegraphics[width=1.0\linewidth]{../figures/Ch3/ch3IWRIk.png}
    \caption{Hydraulic response of biochar-treated soil: (a) detected suction range, (b) change of water retention capacity, (c) penalty of hydraulic conductivity, and (d) relationship between $I_{\mathrm{WR}}$ and $I_{\mathrm{k}}$}
    \label{fig:ch3IWRIk}
\end{figure}

\subsection{Biochar-induced pore structure redistribution}
The change in mercury intrusion volume relative to the untreated specimen for each pore domain can be calculated by:
\begin{equation}
    \Delta V \left( i \right) = V \left( i \right) - V_0 \left( i \right)
\end{equation}
where $V(i)$ is the absolute volume within the domain $i$, and $V_0(i)$ is the volume of the untreated soil at that same \ac{DC}. A comparison of $\Delta V(i)$ is summarised in Figure \ref{fig:ch3Poredomain}. Figure \ref{fig:ch3mechanism} links these structural changes to the suction-dependent hydraulic response.  At 95\% \ac{DC}, the retention pore volume increases by similar amounts in D-B5 and D-B10 (0.041 and 0.039 mL/g), while $D_\mathrm{m}$ increases and $D_\mathrm{L}$ remains unchanged. This indicates that fine-pore development becomes packing-limited at about 5\% biochar. With little free volume, bochar particles and fragments occupy aggregate contacts, creating fine soil–biochar interfacial pores while constricting drainage-domain throats, as supported by the \ac{SEM} observations. Accordingly, D-B5 shows the largest drainage-domain loss, the highest $d_\mathrm{s}$ (2.7943), and an increase in \ac{AEV}. At 10\% biochar, additional soil-biochar and biochar-biochar interfaces form local voids that partly restore moderate-sized pathways, reducing the \ac{AEV} to 5.017 kPa and the unsaturated HC penalty. The dense response therefore reflects non-monotonic competition between throat constriction and interface-pore formation. This interpretation is consistent with evidence that compaction mainly alters inter-aggregate macropores while finer intra-aggregate pores are less sensitive (Sun et al. 2020). At 80\% \ac{DC}, the open inter-aggregate fabric allows biochar to occupy and subdivide pre-existing macropores. Accordingly, Figure \ref{fig:ch3Poredomain} shows increases in both the retention and drainage domains at the expense of the interface domain. $D_\mathrm{s}$ decreases in L-B5 and L-B10 while the \ac{AEV} increases, indicating progressive refinement of the pore throats controlling air entry. Although newly formed interfacial voids partly offset macropore loss at 10\% biochar, subdivision produces more numerous but narrower and more tortuous flow paths. This increases water storage while reducing integrated \ac{HC} in the interface and drainage domains. This increases water storage while reducing integrated \ac{HC} in the interface and drainage domains, consistent with the storage–transmission trade-off associated with pore redistribution and increased tortuosity reported by Hussain and Ravi \parencite*{hussainInvestigatingUnsaturatedHydraulic2021}.\par
Near saturation and at low suction, \ac{HC} is governed mainly by the connectivity and controlling throat size of the interface and drainage pore domains. Once suction exceeds the \ac{AEV}, air progressively invades the larger pores and disrupts the main flow paths, causing \ac{HC} to decrease \parencite{fredlundPredictingPermeabilityFunction1994}. Water is then increasingly retained in finer pores and as menisci or surface films, with the residual fine-pore network becoming more important for water transmission \parencite{tullerHydraulicConductivityVariably2001}. The similar $D_\mathrm{m}$ values of the 5\% and 10\% treatments (2.05-2.15 $\mu$m at both \acp{DC}) are consistent with the convergence of their \ac{HC} curves beyond approximately 60 kPa.\par
Overall, biochar restricts low-suction transmission through pore-throat constriction and increased tortuosity, while enhancing water storage and, in some cases, maintaining residual liquid continuity at higher suction. The hydraulic response therefore depends not only on pore-volume redistribution, but also on pore connectivity, constriction, tortuosity and wettability.\par
\begin{figure}[htbp]
    \centering
    \includegraphics[width=0.6\linewidth]{../figures/Ch3/ch3Poredomain.png}
    \caption{Change of MIP-detected volume relative to the untreated specimen in different pore domains}
    \label{fig:ch3Poredomain}
\end{figure}

\begin{figure}[htbp]
    \centering
    \includegraphics[width=0.8\linewidth]{../figures/Ch3/ch3mechanism.png}
    \caption{Mechanism of compaction-dependent soil-biochar-water interactions}
    \label{fig:ch3mechanism}
\end{figure}

\section{Conclusions}
This chapter established a mechanistic link between initial compaction, biochar dosage, pore restructuring and the suction-dependent hydraulic behaviours of biochar-amended soils. The main conclusions are as follows:\par
Initial compaction determined how biochar is accommodated within the soil fabric. At lower \ac{DC}, biochar mainly occupies and subdivides inter-aggregate macropores, leading to progressive pore refinement and increased water storage. At higher \ac{DC}, limited free volume promotes soil-biochar interfacial pores and drainage-throat constriction. Fine-pore development becomes packing-limited at approximately 5\% biochar, while further addition created interfacial voids that partly restored larger pathways. Biochar-induced pore restructuring is therefore progressive under loose compaction but non-monotonic under dense compaction. The pore domains controlling hydraulic behaviour change with the drying process. At low suction ranges, water storage and hydraulic conductivity are governed mainly by the larger interface and drainage pores. As these pores drain, fine retention pores and surface-water films become increasingly important for both retained water and residual transmission. Pore-volume redistribution therefore explains changes in water storage, whereas unsaturated HC additionally depends on pore connectivity, throat constriction and tortuosity.\par
Enhanced water storage cannot generally lead to enhanced hydraulic transmission. Biochar increases retention most clearly under loose compaction but reduces conductivity within the drainage domain. At higher suction, however, additional fine-pore storage can coexist with improved residual flow after the larger pores have drained. The storage-transmission response to biochar is controlled by the combined effects of compaction state, pore domain and suction. From an engineering perspective, these findings support zone-specific biochar application in slopes, embankments and vegetated earthworks. In densely compacted or deep-layer fills, lower hydraulic conductivity may restrict rapid water migration but may also delay drainage and prolong elevated water contents after infiltration. The limited additional benefit above 5\% biochar should also be considered. In more loosely compacted surface soils, increased retention may improve plant water availability, while reduced drainage-domain conductivity may moderate rapid moisture loss. Biochar dosage should therefore be selected together with the target compaction state and hydraulic function.\par